\documentclass[10pt,a4paper]{amsart}
\usepackage[margin=19mm]{geometry}
\usepackage{amsmath,amssymb,mathtools}
\usepackage[scr=rsfs,cal=euler]{mathalfa}
\usepackage{xfrac}
\usepackage[unicode,hidelinks,bookmarks=false]{hyperref}
\newcommand{\ie}{i.e.\ }
\newcommand{\cf}{cf.\ }
\newcommand{\resp}{respectively\ }
\newcommand{\Subsection}{\subsection}
\providecommand{\nobreakdash}{\nobreak\hbox{-}}
\setlength{\emergencystretch}{2em}
\multlinegap0pt
\usepackage{mathtools}
\usepackage{amssymb}
\usepackage{bm}
\usepackage{dsfont}
\newtheorem{proposition}{Proposition}
\newtheorem{convention}{Convention}
\newtheorem{lemma}{Lemma}
\newtheorem{remark}{Remark}
\newcommand{\Diff}{\mathrm{Diff}(M)}
\newcommand{\Forms}[1]{\Omega^{#1}(M)}
\newcommand{\GenTM}{\mathbb TM}
\newcommand{\Id}{\mathrm{Id}}
\newcommand{\LieD}{\mathcal{L}}
\newcommand{\N}{\mathbb{N}}
\newcommand{\Ucal}{\mathcal{U}}
\newcommand{\VFs}{\mathfrak{X}(M)}
\renewcommand{\d}{\mathrm{d}}
\newcommand{\eh}[1]{\mathrm{exp}^h(#1)}
\newcommand{\exh}{\mathrm{exp}^h}
\newcommand{\tens}{\otimes}
\newcommand{\tensf}[2]{\mathcal{T}^{#1}_{#2}(M)}
\title{The group of autoequivalences of an exact Courant algebroid as a tame Fréchet Lie group}
\author{Jan Niklas Heck}
\date{Source: arXiv:2608.16821v1 (CC BY 4.0)}
\begin{document}
\maketitle

\noindent\emph{Source and licence.} The group of autoequivalences of an exact Courant algebroid as a tame Fréchet Lie group. Version arXiv:2608.16821v1 (CC BY 4.0), distributed by arXiv under the Creative Commons Attribution 4.0 International licence, http://creativecommons.org/licenses/by/4.0/, as recorded in the arXiv licence metadata for this exact version and shown on the abstract page https://arxiv.org/abs/2608.16821v1. Full licence notice: \texttt{licenses/arxiv-2608.16821-CC-BY-4.0.txt}.\par\smallskip

\noindent\textit{Editorial scope.} Counted unit: the article's lemma lemma (source0.tex lines 3020--3032). The article's complete original proof of it is delivered with it (source0.tex lines 3033--3053, one \texttt{proof} environment of 21 lines). No mathematics is written, repaired or re-derived: the statement and the proof are the article's own lines, in the article's own notation, and no retained line is substituted or re-flowed.  Retained uncounted as the source-local context the proof needs: lines 1120--1123; lines 2097--2100; lines 2267--2270; lines 2478--2482; lines 2484--2488; lines 2508--2521.  Nothing was omitted from any retained span, so the package carries no omission record.  No retained line contains a citation, so the wrapper carries no bibliography and no bibliography entry.  No retained line contains a figure environment, an image-inclusion command or drawing code, so no image is delivered and the package declares no asset dependency.  Editorial formatting only, in addition: an AMS \texttt{amsart} preamble that restates the article's own declarations and macros used by the retained lines, namely the environments \texttt{proposition}, \texttt{convention}, \texttt{lemma}, \texttt{remark} on separate counters, together with the standard packages those lines need.  The retained environments print in this document as Proposition 1, Convention 1, Lemma 1, Remark 1 in order of appearance, whereas the article prints the same items with the numbering of its own full document.  The complete native source of the article as distributed in the arXiv e-print for this exact version is bundled unchanged as \texttt{sources/207858/source0.tex}.
\begin{proposition}
\label{Prop:EplusFsect}
Let $E\to M$ and $F\to M$ be vector bundles and $M$ a compact manifold. Then $\Gamma(E\oplus F)\cong\Gamma(E)\times\Gamma(F)$ is a zero-tame isomorphism.
\end{proposition}

\begin{convention}
\label{Conv:ChartsOnDiff}
Using the notation of the previous remark we denote by $(\Ucal,\exh_f)$, the local parametrization of $\Diff$ around $f\in \Diff$ defined by the diffeomorphism $\exh_f\coloneqq R_f\circ\exh:\Ucal\to R_f(\eh{\Ucal})\subseteq \Diff$. Note that its inverse is a chart around $f\in \Diff$. For convenience, we denote $(\Ucal,\exh_{\Id})$ by $(\Ucal,\exh)$.
\end{convention}

\begin{convention}
\label{Conv:tensfields}
As there are different conventions in literature, we want to clarify that in the following we will use the notation $\tensf{n}{m}\coloneqq \Gamma((TM)^{\tens n}\tens(T^*M)^{\tens m})$, $n,m\in \N$, to denote the smooth $(n,m)$-tensor fields on a (compact) manifold $M$.
\end{convention}

\begin{convention}
\label{Conv:StandChartOnDifftimesTensOrForms}
Let $(\Ucal,\exh_f)$ be the standard local parametrization of $\Diff$ centered at $f\in \Diff$ (cf.~Convention \ref{Conv:ChartsOnDiff}). We then denote the resulting standard local parametrization of $\Diff \times \tensf{n}{m}$, $n,m \in \N$, centered at $(f,0)$ by
$\exh_{f,\mathcal T}\coloneqq \exh_f\times \mathrm{Id}_{\tensf{n}{m}}$ and similarly we denote the standard local parametrization of $\Diff \times \Forms{k}$, $k\in \N$, centered at $(f,0)$ by $\exh_{f,\Omega}\coloneqq \exh_f \times \mathrm{Id}_{\Forms{k}}$. Note that $\exh_{f,\mathcal T}$, and $\exh_{f,\Omega}$ are tame diffeomorphisms onto their image and so the inverses of these maps are local charts of the respective manifolds. For $f=\Id$ we will again denote the respective local parametrizations by $\exh_{\mathcal{T}}\coloneqq \exh_{\Id,\mathcal{T}}$ and $\exh_{\Omega}\coloneqq \exh_{\Id,\Omega}$. The respective degrees will be clear from context. 
\end{convention}

\begin{lemma}
\label{Lemma:PullbackDifferential}
Let $P:\Diff\times \tensf{n}{m}\to \tensf{n}{m}$, $n,m\in \N$, be the pullback as in the previous theorem. Then in the chart $\exh_{f,\mathcal T}$ (cf.~Convention \ref{Conv:StandChartOnDifftimesTensOrForms}) centered at $(f,0)\in \Diff\times \tensf{n}{m}$ we find $D(P\circ \exh_{f,\mathcal T})(0,\tau)(X,T)=f^*\LieD_X \tau+f^*T$ for
$\tau,T\in \tensf{n}{m}$, and $X\in \VFs$.  
\end{lemma}

\begin{remark}
\label{Rem:PushForwDifferential}
Let $V$ denote the inversion map of $\Diff$. The pushforward of a tensor field is the map $P\circ (V\times \mathrm{Id}_{\tensf{n}{m}}):\Diff\times \tensf{n}{m}\to \tensf{n}{m}$, $n,m\in \N$, which is smooth tame as a composition of smooth tame maps. 
Completely analogously to Lemma \ref{Lemma:PullbackDifferential} we can now take the flow $\Phi_t^X$ of a vector field $X\in \VFs$ and calculate for $\tau \in \tensf{n}{m}$ the differential
\begin{align*}
    \left.\frac{\d}{\d t}\right|_{t=0}P(V(\eh{tX}),\tau)=\left.\frac{\d}{\d t}\right|_{t=0}P(V(\Phi_t^X),\tau)=\left.\frac{\d}{\d t}\right|_{t=0}P(\Phi_{-t}^X,\tau)=-\LieD_X\tau.
\end{align*}
As before this shows that in the standard chart $\exh_{f,\mathcal T}$ (cf.~Convention \ref{Conv:StandChartOnDifftimesTensOrForms}) around $(f,0)\in \Diff\times \tensf{n}{m}$ we find
\begin{align*}
    D(P\circ (V\times \mathrm{Id}_{\Forms{n}})\circ \exh_{f,\mathcal T})(0,\tau)(X,T)&=\left(\left.\frac{\d}{\d t}\right|_{t=0}P(\Phi_{-t}^X,f_*\tau)\right)+f_*T\\
    &=-\LieD_X(f_*\tau)+f_*T
\end{align*}
for $X\in \VFs$ and $T\in \tensf{n}{m}$.
\end{remark}

\begin{lemma}
Let $M$ be a compact manifold. Then the map 
\begin{align*}
    c_\Omega:\Forms{2}\times \Gamma(\mathrm{End}(\GenTM))&\longrightarrow \Gamma(\mathrm{End}(\GenTM))\\
    (B,F)&\longmapsto e^B\circ F\circ e^{-B} 
\end{align*}
and the map
\begin{align*}
    c_D:\Diff\times \Gamma(\mathrm{End}(\GenTM))&\longrightarrow \Gamma(\mathrm{End}(\GenTM))\\
    (f,F)&\longmapsto f_*\circ F\circ f_*^{-1}
\end{align*}
are smooth tame.
\end{lemma}

\begin{proof}
The fact that the map $c_\Omega$ is smooth tame follows directly from it being a nonlinear vector bundle operator up to the zero-tame identification 
$\Forms{2}\times \Gamma(\mathrm{End}(\GenTM))\cong \Gamma((\bigwedge^2T^*M)\oplus\mathrm{End}(\GenTM))$ (cf.~Proposition \ref{Prop:EplusFsect}).

Let us now consider $c_D$. We note that $\mathrm{End}(\GenTM)\cong (TM\oplus T^*M)\tens (TM\oplus T^*M)^*$. Again by Proposition \ref{Prop:EplusFsect}, this induces a zero-tame isomorphism
$\Gamma(\mathrm{End}(\GenTM))\cong \tensf{1}{1}\times \tensf{2}{0}\times \tensf{0}{2}\times \tensf{1}{1}$ with $\tensf{n}{m}$ as in Convention \ref{Conv:tensfields}. So, it suffices to check how $c_D$ acts on these components. Let $J\in \Gamma(\mathrm{End}(E))$ correspond to
\begin{align*}
    \begin{pmatrix}
        J_1 & J_2\\
        J_3 & J_4
\end{pmatrix}
\end{align*}
with $J_1\in \tensf{1}{1}$, $J_2\in \tensf{2}{0}$, $J_3\in \tensf{0}{2}$, and $J_4\in \tensf{1}{1}$ then a short calculation shows that $c_D(f,J)$ corresponds simply to
\begin{align*}
    \begin{pmatrix}
        f_*J_1 & f_*J_2\\
        f_*J_3 & f_*J_4
\end{pmatrix}
\end{align*}
for every $f\in \Diff$. Hence, the map $c_D$ is smooth tame as the push forward of tensor fields along diffeomorphisms is a smooth tame map (cf.~Remark \ref{Rem:PushForwDifferential}).
\end{proof}

\end{document}
